\documentclass[12pt,oneside,openany]{book}
% ============================================================
%  Preamble -- Quantum Theory (F. P. Schuller); machinery from the Brunton/Mehrling notes
% ============================================================
\usepackage[T1]{fontenc}
\usepackage[utf8]{inputenc}
\usepackage{amsmath,amsthm,mathtools}
\usepackage{newpxtext,newpxmath} % provides the AMS symbols: do not load amssymb
\usepackage{bm}
\usepackage{microtype}
% Screen layout (31.5x18 cm = Acer Aspire E15 screen minus 15 mm (left/right) and 7 mm (top/bottom) viewer margin, PDF viewer in reading mode, not printed): two parallel columns (paracol), the text column
% 18.8 cm and a 10 cm right column (was 15 + 6.5 until 2026-11-05) for notes, answers, feedback,
% examples, proofs, course notes, history and asides. Every right-column item starts level with the
% paragraph it belongs to; if the right column is still busy, the text waits (white space).
% The old A4 print layout was: a4paper,inner=2.4cm,outer=3.6cm,marginparwidth=3.0cm,marginparsep=0.3cm (twoside).
\ifdefined\scrollmode % continuous-scroll build (scroll.py): tall header-free pages, stacked afterwards
  \usepackage[paperwidth=31.5cm,paperheight=100cm,top=0.3cm,bottom=0.3cm,left=0.9cm,textwidth=29.3cm,
    marginparsep=0.5cm,marginparwidth=10cm]{geometry}
\else
\usepackage[paperwidth=31.5cm,paperheight=18cm,includehead,top=0.4cm,bottom=1.5cm,left=0.9cm,textwidth=29.3cm,
  marginparsep=0.5cm,marginparwidth=10cm,headheight=14pt,headsep=0.15cm]{geometry} % one-line header holds the page number: no footer
\fi
\usepackage{paracol}
\usepackage{needspace} % \Needspace*{10.5cm} before a tall unbreakable right-column box (history) that would run off the page
\setcolumnwidth{18.8cm/0.5cm,10cm}
% counters shared by the two columns (equations in a proof, examples, answers ... in the right column)
% (sectioning counters cannot be global with this kernel: they are copied to the right column by
% \synccounter before each switch, see \side@emit)
\globalcounter{equation}\globalcounter{footnote}\globalcounter{figure}
\usepackage{booktabs,array,tabularx,longtable}
\usepackage[dvipsnames]{xcolor}
\usepackage[shortlabels,inline]{enumitem}
\usepackage{tikz}
\usetikzlibrary{arrows.meta,positioning,decorations.pathreplacing,calc}
\usepackage{pgfplots}
\pgfplotsset{compat=1.18}
\usepackage[most]{tcolorbox}
\usepackage{fancyhdr}
\usepackage[colorlinks=true,linkcolor=NavyBlue,urlcolor=NavyBlue,citecolor=NavyBlue]{hyperref}
\usepackage[nameinlink]{cleveref}
\usepackage{listings}

\setlist{itemsep=2pt,topsep=4pt}
\setlength{\parskip}{3pt}
\setlength{\emergencystretch}{2em}
\hfuzz=5pt

% ---------- headers ----------
\pagestyle{fancy}
\fancyhf{}
\fancyhead[L]{\small\leftmark}
\fancyhead[R]{\small\makebox[\marginparwidth][r]{\rightmark\quad\thepage}}
\renewcommand{\headrulewidth}{0.3pt}
\fancypagestyle{plain}{\fancyhf{}\fancyhead[R]{\small\makebox[\marginparwidth][r]{\thepage}}\renewcommand{\headrulewidth}{0pt}}
\ifdefined\scrollmode
  \raggedbottom \pagestyle{empty}\fancypagestyle{plain}{\fancyhf{}\renewcommand{\headrulewidth}{0pt}}
  % first and last page of the \include d chapter, for scroll.py
  \newcounter{scrollabs}\newwrite\scrollpg\immediate\openout\scrollpg=\jobname.scrollpages
  \AddToHook{shipout/before}{\stepcounter{scrollabs}}
  \AddToHook{include/before}{\immediate\write\scrollpg{\the\numexpr\value{scrollabs}+1\relax}}
  \AddToHook{include/after}{\immediate\write\scrollpg{\the\value{scrollabs}}}
  % a chapter longer than ~3 pages of 100 cm starts a new long page at the next \section (clean seam, PDF viewers
  % stop at ~508 cm); the page number of every break goes to the same file, scroll.py stacks accordingly
  \newcounter{scrollsince}
  \newcommand{\scrollbreak}{\clearpage\immediate\write\scrollpg{\the\numexpr\value{scrollabs}+1\relax}\setcounter{scrollsince}{\value{scrollabs}}}
  \AddToHook{include/before}{\setcounter{scrollsince}{\value{scrollabs}}}
  \AddToHook{cmd/section/before}{\ifnum\numexpr\value{scrollabs}-\value{scrollsince}\relax>2 \scrollbreak\fi}
\fi

% ---------- colours (bright on purpose: yellow notes, blue definitions, red theorems, green examples) ----------
\definecolor{defcol}{HTML}{1565C0}  % blue: definitions
\definecolor{thmcol}{HTML}{D32F2F}  % red: theorems, propositions, lemmas, corollaries
\definecolor{intcol}{HTML}{2E7D32}  % green: examples (also the green of the figures)
\definecolor{notecol}{HTML}{FBC02D} % yellow: margin notes and their answers
\definecolor{intucol}{HTML}{8E24AA} % purple: intuition
\definecolor{fincol}{HTML}{EF6C00}  % orange: applications
\definecolor{cscol}{HTML}{00897B}   % teal: case studies
\definecolor{keycol}{HTML}{C2185B}  % magenta: key formulas
\definecolor{exercol}{HTML}{3949AB} % indigo: exercises
\definecolor{srccol}{HTML}{546E7A}  % grey-blue: sources, course notes, remarks
\definecolor{histcol}{HTML}{B8860B} % bronze: history boxes (classical columns)
\definecolor{recallcol}{HTML}{00ACC1} % cyan: recall boxes (material already seen earlier)
\definecolor{doubtcol}{HTML}{C62828}
\definecolor{ideacol}{HTML}{1565C0}
\definecolor{okcol}{HTML}{388E3C}
\colorlet{anscol}{notecol!85!black}

% ---------- table rows: a little air between lines in every table ----------
\AddToHook{env/tabular/begin}{\renewcommand{\arraystretch}{1.35}}
\AddToHook{env/tabularx/begin}{\renewcommand{\arraystretch}{1.35}}
\AddToHook{env/longtable/begin}{\renewcommand{\arraystretch}{1.35}}

% ---------- theorem-like environments (numbered by chapter), coloured headings ----------
\newtheoremstyle{coldef}{}{}{\normalfont}{}{\bfseries\color{defcol}}{.}{ }{}
\newtheoremstyle{colex}{}{}{\normalfont}{}{\bfseries\color{intcol}}{.}{ }{}
\newtheoremstyle{colexer}{}{}{\normalfont}{}{\bfseries\color{exercol}}{.}{ }{}
\newtheoremstyle{colthm}{}{}{\itshape}{}{\bfseries\color{thmcol}}{.}{ }{}
\newtheoremstyle{colrem}{}{}{\normalfont}{}{\itshape\color{srccol}}{.}{ }{}
\theoremstyle{coldef}
\newtheorem{definition}{Definition}[chapter]
\theoremstyle{colex}
\newtheorem{example}[definition]{Example}
\theoremstyle{colexer}
\newtheorem{exercise}{Exercise}[chapter]
\theoremstyle{colthm}
\newtheorem{theorem}[definition]{Theorem}
\newtheorem{proposition}[definition]{Proposition}
\newtheorem{lemma}[definition]{Lemma}
\newtheorem{corollary}[definition]{Corollary}
\theoremstyle{colrem}
\newtheorem{remark}[definition]{Remark}
\crefname{exercise}{Exercise}{Exercises}
\globalcounter{definition}\globalcounter{exercise}

% Right-column items are queued and written into the right column at the next point where the
% text column is between paragraphs at top level (end of the paragraph, end of a box, end of a
% list): there the columns are synchronised, the queued items are typeset in the right column,
% and the text continues from where it was.
\makeatletter
\newif\ifnote@inbox
\newif\ifside@inright
\def\note@queue{}
\def\side@pcname{paracol}
\newcommand{\side@emit}{\ifx\note@queue\@empty\else
  \let\note@tmp\note@queue\gdef\note@queue{}%
  \synccounter{chapter}\synccounter{section}\synccounter{subsection}\synccounter{subsubsection}%
  \switchcolumn*\side@inrighttrue\note@tmp\par\side@inrightfalse\switchcolumn\fi}
\newcommand{\note@flush}{\ifvmode\ifnote@inbox\else\ifside@inright\else
  \ifx\@currenvir\side@pcname\side@emit\fi\fi\fi\fi}
\newcommand{\side@put}[1]{\gappto\note@queue{#1}\note@flush}
\AddToHook{para/after}{\note@flush}
\tcbset{note queue/.code={\appto\kvtcb@afterbox{\note@flush}}}
% every chapter/appendix file is one paracol environment: it starts right after the \chapter
% line of the file (\begin{paracol}{2}) and is closed here (appendix A, only long tables, has none:
% longtable does not work inside paracol)
\AddToHook{include/end}{\par\note@flush\ifx\@currenvir\side@pcname\end{paracol}\fi}
\newcommand{\flushside}{\par\note@flush}
\makeatother
\tcbset{thmbox/.style={enhanced,breakable,colback=#1!7,colframe=#1!7,borderline west={3pt}{0pt}{#1},
  sharp corners,left=6pt,right=4pt,top=2pt,bottom=2pt,before skip=8pt,after skip=8pt,
  before upper={\csname note@inboxtrue\endcsname},note queue}}
\tcolorboxenvironment{definition}{thmbox=defcol}
\tcolorboxenvironment{theorem}{thmbox=thmcol}
\tcolorboxenvironment{proposition}{thmbox=thmcol}
\tcolorboxenvironment{lemma}{thmbox=thmcol}
\tcolorboxenvironment{corollary}{thmbox=thmcol}
% (examples are boxed in the right column, see \begin{example} below)
\tcolorboxenvironment{exercise}{thmbox=exercol}
\tcolorboxenvironment{remark}{thmbox=srccol}

% ---------- content boxes ----------
\tcbset{notebox/.style={enhanced,breakable,colback=#1!7,colframe=#1,
  fonttitle=\bfseries,coltitle=white,colbacktitle=#1,boxrule=0.6pt,arc=2pt,
  left=6pt,right=6pt,top=4pt,bottom=4pt,before skip=10pt,after skip=10pt,
  before upper={\csname note@inboxtrue\endcsname},note queue}}
% Box title: the box type in capitals, then the topic
\newcommand{\boxtitle}[2]{\MakeUppercase{#1}\ifx&#2&\else: #2\fi}
% Intuition: the "physicist's reading" of a result
\newtcolorbox{intuition}[1][]{notebox=intucol,title={\boxtitle{Intuition}{#1}}}
% Application: where a method is used in science and engineering (orange)
\newtcolorbox{application}[1][]{notebox=fincol,title={\boxtitle{Application}{#1}}}
% Case study: a data set worked through in the lectures (cylinder wake, Lorenz, measles, ...)
\newtcolorbox{casestudy}[1][]{notebox=cscol,title={\boxtitle{Case study}{#1}}}
% Summary / key formulas
\newtcolorbox{keyformulas}{notebox=keycol,fontupper=\small,title={\boxtitle{Key formulas}{}}}
\newtcolorbox{keyconcepts}{notebox=keycol,fontupper=\small,title={\boxtitle{Key concepts}{}}}
% Sources: which lectures, videos and files a chapter is built on
\newtcolorbox{sources}{notebox=srccol,fontupper=\footnotesize,title={\boxtitle{Sources for this chapter}{}},before upper=\raggedright}
\newcolumntype{L}{>{\raggedright\arraybackslash}X}
\newcolumntype{V}{>{\raggedright\arraybackslash}p{4.6cm}} % video titles in the Sources box
% Reading: pointers to the other Schuller notes sets and to his papers
\newtcolorbox{reading}[1][]{notebox=srccol,title={\boxtitle{Reading}{#1}}}

% non-floating figure (floats inside paracol columns can leave a page almost empty): same caption
% and numbering as figure
\makeatletter
\newenvironment{nfigure}{\par\medskip\begin{center}\def\@captype{figure}}{\end{center}\medskip}
\makeatother
% Note on course material (errata, missing lectures, differences 2013/2024): right column, see below
% Secondary material ("beyond the core"): can be skipped on a first reading
\newtcolorbox{secondary}[1][]{enhanced,breakable,colback=white,colframe=white,boxrule=0pt,
  borderline west={1.5pt}{0pt}{srccol!60},sharp corners,left=8pt,right=0pt,top=1pt,bottom=1pt,
  before skip=6pt,after skip=6pt,fontupper=\small,
  title={\footnotesize\sffamily$\diamond$\ \boxtitle{Beyond the core}{#1}},
  coltitle=srccol,colbacktitle=white,fonttitle=\bfseries,titlerule=0pt,toptitle=0pt,bottomtitle=0pt,
  before upper={\csname note@inboxtrue\endcsname},note queue}
% ============================================================
%  The right column: the reader's notes/answers/feedback in \scriptsize; examples, proofs, course notes,
%  history and asides in \footnotesize
% ============================================================
%  \note{...}     the reader's comment/doubt, NEW (not yet read by Claude)
%  \note*{...}    a request that has been carried out (green DONE tag)
%  \note+{...}    a comment that has been read and needs no answer (blue SEEN tag)
%  \begin{answer}{question}[topic] ... \end{answer}
%                 a doubt and its answer in one box: the question on top, then a dark yellow
%                 bar "ANSWER n: topic" and the answer; \answerpart{topic} starts a further
%                 answer in the same box (one per question); \label{ans:...} makes it citable
%  \feedback{...} a comment on one of the reader's own solutions (indigo)
%  \begin{history}[topic], \begin{aside}[topic], \begin{coursenote}[topic]
%  example and proof environments are also printed here (sideexample = example)
%  Inside a coloured box of the main text, margin items appear right after the box.
\makeatletter
\tcbset{sidebox/.style={enhanced,colback=#1!25,colframe=#1,boxrule=0.5pt,arc=1pt,
    left=3pt,right=3pt,top=2pt,bottom=2pt,before skip=0pt,after skip=0pt,
    fontupper=\scriptsize\sffamily\mdseries\upshape\raggedright,fonttitle=\scriptsize\sffamily\bfseries,
    before upper={\note@inboxtrue}}}
\newcommand{\side@tag}[2]{\colorbox{#1}{\textcolor{white}{\bfseries\scriptsize #2}}\par}
% Notes and answered notes are numbered per chapter like footnotes: a small dark-yellow number in
% the text marks the point where the note was written, the same number starts its box.
\newcounter{sidenote}[chapter]
\globalcounter{sidenote}
\newcommand{\side@numput}[1]{\stepcounter{sidenote}%
  \ifhmode\ifdim\lastkern=0.01pt\textsuperscript{\sffamily\bfseries\color{anscol!75!black},}\fi
    \textsuperscript{\sffamily\bfseries\color{anscol!75!black}\thesidenote}\kern0.01pt\fi
  \begingroup\edef\x{\endgroup\noexpand\side@put{\def\noexpand\side@num{\thesidenote}\unexpanded{#1}}}\x}
\newcommand{\side@numtag}{{\bfseries\color{anscol!75!black}\side@num}\enspace}
\NewDocumentCommand{\note}{s t+ +m}{\side@numput{\begin{tcolorbox}[sidebox=notecol]%
  \IfBooleanT{#1}{\side@tag{okcol}{DONE}}\IfBooleanT{#2}{\side@tag{defcol}{SEEN}}\side@numtag#3\end{tcolorbox}}}
\newcommand{\feedback}[1]{\side@put{\begin{tcolorbox}[sidebox=exercol]\side@tag{exercol}{FEEDBACK}#1\end{tcolorbox}}}
\newcounter{answer}[chapter]
\renewcommand{\theanswer}{\thechapter.\arabic{answer}}
\crefname{answer}{answer}{answers}
\globalcounter{answer}
\newcommand{\answerpart}[1]{\refstepcounter{answer}%
  \tcbsubtitle[colback=anscol,colframe=anscol,coltext=black,fontupper=\scriptsize\sffamily\bfseries,
    left=3pt,top=0.5pt,bottom=0.5pt,before skip=2pt,after skip=2pt]{ANSWER \theanswer\ifx&#1&\else: #1\fi}}
\NewDocumentEnvironment{answer}{m O{} +b}{\side@numput{\begin{tcolorbox}[sidebox=notecol,colback=notecol!12]%
  \side@numtag#1\answerpart{#2}#3\end{tcolorbox}}}{}
% Content boxes of the right column (examples, proofs, course notes, history, asides): \footnotesize,
% one step above the \scriptsize of the reader's notes
\tcbset{sidecontent/.style={sidebox=#1,colback=#1!6,breakable,
    fontupper=\footnotesize\mdseries\upshape\raggedright,fonttitle=\footnotesize\sffamily\bfseries,coltitle=white,colbacktitle=#1}}
\NewDocumentEnvironment{aside}{O{} +b}{\side@put{\begin{tcolorbox}[sidecontent=histcol,
  title={\boxtitle{Aside}{#1}}]#2\end{tcolorbox}}}{}
% Recall: a repetition of something explained earlier (another chapter, earlier in this one),
% restated briefly with a \cref to the original, so that it is visibly a repetition.
% Longer repeated passages go verbatim to a "Set aside" section at the end of the chapter.
\NewDocumentEnvironment{recall}{O{} +b}{\side@put{\begin{tcolorbox}[sidecontent=recallcol,
  title={\boxtitle{Recall}{#1}}]#2\end{tcolorbox}}}{}
\NewDocumentEnvironment{coursenote}{O{} +b}{\side@put{\begin{tcolorbox}[sidecontent=srccol,
  title={\boxtitle{Course note}{#1}}]#2\end{tcolorbox}}}{}
% Examples and proofs keep their theorem-like look (numbering, labels, QED) but go to the right column
\let\origexample\example \let\endorigexample\endexample
\RenewDocumentEnvironment{example}{o +b}{\side@put{\begin{tcolorbox}[thmbox=intcol,
  fontupper=\footnotesize,before skip=0pt,after skip=0pt]%
  \IfValueTF{#1}{\begin{origexample}[#1]}{\begin{origexample}}#2\end{origexample}\end{tcolorbox}}}{}
\NewDocumentEnvironment{sideexample}{o +b}{\IfValueTF{#1}{\begin{example}[#1]}{\begin{example}}#2\end{example}}{}
\let\origproof\proof \let\endorigproof\endproof
% A proof must start level with the top of its theorem: a theorem directly followed by \begin{proof}
% is held back (\thm@pend) and typeset by the proof, right after the proof has been queued.
\def\thm@pend{}
\newtoks\thm@tk
\newcommand{\thm@typeset}{{\edef\@currenvir{\side@pcname}\begin{tcolorbox}[thmbox=thmcol]\expandafter\IfValueTF\expandafter{\thm@opt}%
  {\begin{orig\thm@cur}[\thm@opt]\the\thm@tk\end{orig\thm@cur}}{\begin{orig\thm@cur}\the\thm@tk\end{orig\thm@cur}}\end{tcolorbox}}}
\newcommand{\thm@peek}[3]{\def\thm@a{proof}\def\thm@b{#3}%
  \ifx\thm@a\thm@b\gdef\thm@pend{\thm@typeset}\else\thm@typeset\fi#1#2{#3}}
\newcommand{\thm@look}[2]{\@ifnextchar\begin{\thm@peek{#1{#2}}}{\thm@typeset#1{#2}}}
\RenewDocumentEnvironment{proof}{o +b}{\par{\edef\@currenvir{\side@pcname}\side@put{\begin{tcolorbox}[thmbox=thmcol,colback=thmcol!3,
  fontupper=\footnotesize,before skip=0pt,after skip=0pt]%
  \IfValueTF{#1}{\begin{origproof}[#1]}{\begin{origproof}}#2\end{origproof}\end{tcolorbox}}%
  \thm@pend}\gdef\thm@pend{}}{}
\@for\thm@n:=theorem,proposition,lemma,corollary\do{%
  \RemoveFromHook{env/\thm@n/before}[tcolorbox]\RemoveFromHook{env/\thm@n/after}[tcolorbox]%
  \expandafter\let\csname orig\thm@n\expandafter\endcsname\csname\thm@n\endcsname
  \expandafter\let\csname endorig\thm@n\expandafter\endcsname\csname end\thm@n\endcsname
  \expandafter\edef\csname thm@def@\thm@n\endcsname{\noexpand\RenewDocumentEnvironment{\thm@n}{o +b}%
    {\noexpand\global\noexpand\thm@tk{##2}\noexpand\gdef\noexpand\thm@cur{\thm@n}\noexpand\gdef\noexpand\thm@opt{##1}\noexpand\thm@look}{}}%
  \csname thm@def@\thm@n\endcsname}
\makeatother
% History: a fact from the history of finance/mathematics, framed like a classical
% temple front: a dentilled entablature (the title) on two Ionic columns (not breakable)
\newcommand{\histcolumn}[2]{% #1 = foot of the column, #2 = top of the capital (coordinate names)
  \fill[histcol] ([xshift=-6pt]#1) rectangle ([xshift=6pt,yshift=2.5pt]#1);
  \fill[histcol] ([xshift=-4.5pt,yshift=2.5pt]#1) rectangle ([xshift=4.5pt,yshift=4pt]#1);
  \fill[histcol!15] ([xshift=-3.5pt,yshift=4pt]#1) rectangle ([xshift=3.5pt,yshift=-5pt]#2);
  \foreach \dx in {-3.5,-1.2,1.2,3.5}
    \draw[histcol,line width=0.4pt] ([xshift=\dx pt,yshift=4pt]#1) -- ([xshift=\dx pt,yshift=-5pt]#2);
  \fill[histcol] ([xshift=-4.5pt,yshift=-5pt]#2) rectangle ([xshift=4.5pt,yshift=-2.5pt]#2);
  \foreach \dx in {-4.5,4.5} {
    \draw[histcol,line width=0.8pt,fill=white] ([xshift=\dx pt,yshift=-4.5pt]#2) circle (1.9pt);
    \fill[histcol] ([xshift=\dx pt,yshift=-4.5pt]#2) circle (0.6pt);}
  \fill[histcol] ([xshift=-6.5pt,yshift=-2.5pt]#2) rectangle ([xshift=6.5pt]#2);}
\makeatletter
\NewDocumentEnvironment{history}{O{} +b}{\side@put{\begin{tcolorbox}[histbox,title={\boxtitle{History}{#1}}]#2\end{tcolorbox}}}{}
\makeatother
\tcbset{histbox/.style={enhanced,colback=histcol!6,colframe=histcol,boxrule=0.8pt,sharp corners,
  fonttitle=\footnotesize\sffamily\bfseries,fontupper=\normalfont\footnotesize\raggedright,coltitle=white,colbacktitle=histcol,
  left=18pt,right=18pt,top=6pt,bottom=3pt,before skip=0pt,after skip=0pt,
  before upper={\csname note@inboxtrue\endcsname},
  overlay={
    \draw[histcol,line width=2.5pt,dash pattern=on 2.5pt off 2pt]
      ([yshift=-1.25pt]title.south west) -- ([yshift=-1.25pt]title.south east);
    \coordinate (hfL) at ([xshift=9pt]frame.south west);
    \coordinate (htL) at ([xshift=9pt,yshift=-2.5pt]title.south west);
    \coordinate (hfR) at ([xshift=-9pt]frame.south east);
    \coordinate (htR) at ([xshift=-9pt,yshift=-2.5pt]title.south east);
    \histcolumn{hfL}{htL}\histcolumn{hfR}{htR}}}}
% The reader's own solution to an exercise, placed right below the exercise
%   \begin{mysolution}[optional topic] ... \end{mysolution}
\newtcolorbox{mysolution}[1][]{notebox=exercol,colback=exercol!5,fontupper=\small,
  title={\boxtitle{My solution}{#1}}}
% A request about the notes themselves (collected in the "Requests" chapter)
%   \request{done|open|noted}{where it was written}{the request}{what was done}
\makeatletter
\newcommand{\request}[4]{\ifstrequal{#1}{done}{\def\req@col{okcol}}{\ifstrequal{#1}{open}{\def\req@col{thmcol}}{\def\req@col{defcol}}}%
  \begin{tcolorbox}[sidebox=notecol,colback=notecol!12,fontupper=\small\sffamily\raggedright,
  before skip=5pt,after skip=5pt]%
  \side@tag{\req@col}{\MakeUppercase{#1}}%
  {\itshape #2.}\ #3%
  \tcbsubtitle[colback=anscol,colframe=anscol,coltext=black,fontupper=\footnotesize\sffamily\bfseries,
    left=3pt,top=0.5pt,bottom=0.5pt,before skip=2pt,after skip=2pt]{WHAT WAS DONE}#4\end{tcolorbox}}
\makeatother



% ---------- video timestamps ----------
%  \ts{4}{1}{3:10}  -> lecture 4 (one video per lecture: the second argument is always 1), at 3:10
%  Each tag is a link to the YouTube video at that second. The IDs (playlist of the 28 lectures)
%  are in videos.tex, generated by ../tools/transcripts.py;
%  an unknown video or a non-numeric time prints the tag without link.
\ExplSyntaxOn
\cs_new_protected:Npn \tsvideo #1#2 { \tl_const:cx { c__ts_#1_tl } { \tl_to_str:n {#2} } }
\int_new:N \l__ts_int
\seq_new:N \l__ts_seq
\cs_new_protected:Npn \tslink #1#2#3#4
  {
    \tl_if_exist:cTF { c__ts_#1@#2_tl }
      { \regex_match:nnTF { \A\d+(:\d+)*\Z } {#3} { \__ts_href:nnnn {#1} {#2} {#3} {#4} } {#4} }
      {#4}
  }
\cs_new_protected:Npn \__ts_href:nnnn #1#2#3#4
  {
    \int_zero:N \l__ts_int
    \exp_args:NNV \seq_set_split:Nnn \l__ts_seq \c_colon_str {#3}
    \seq_map_inline:Nn \l__ts_seq { \int_set:Nn \l__ts_int { 60*\l__ts_int + ##1 } }
    \exp_args:Ne \href { \tl_to_str:n { https://youtu.be/ } \tl_use:c { c__ts_#1@#2_tl } \tl_to_str:n { ?t= } \int_use:N \l__ts_int } {#4}
  }
\ExplSyntaxOff
% Generated by tools/transcripts.py: YouTube ID of every video (chapter@video)
\tsvideo{1@1}{wEI5Mm4P21M} % Experimentum crucis und Quantenwahrscheinlichkeiten
\tsvideo{2@1}{g8xNQl5Rtto} % Pfadintegral und Elementaramplituden
\tsvideo{3@1}{zWzZgNRNpVI} % Wellenfunktionen und Schrödingergleichung
\tsvideo{4@1}{5bC1DvINw3Y} % Einfachspalt und Unschärferelation
\tsvideo{5@1}{YAYEYbae-KU} % Hilbert(t)räume
\tsvideo{6@1}{Wa-Z3NrnEnc} % Spektralsatz (Kleinkinderversion)
\tsvideo{7@1}{YNpeydkLAG4} % Spektralsatz (allgemein)
\tsvideo{8@1}{9ucJEYVFLHo} % Aufgetakelte Hilberträume
\tsvideo{9@1}{bVpeklBliPA} % Allgemeines Rahmenwerk der Quantenmechanik: Zustände und Observablen
\tsvideo{10@1}{5XRHh4bOcOs} % Allgemeines Rahmenwerk der QM: Bestmögliche Vorhersage von Messergebnissen
\tsvideo{11@1}{QunQCGVlT1k} % Allgemeines Rahmenwerk der QM: Dynamik bei und zwischen Messungen
\tsvideo{21@1}{GbqA9Xn_iM0} % Axioms of Quantum Mechanics - Lec01 - Frederic Schuller
\tsvideo{22@1}{Px1Zd--fgic} % Banach Spaces - Lec02 - Frederic Schuller
\tsvideo{23@1}{5OGCIwx0phI} % Separable Hilbert spaces - L03 - Frederic Schuller
\tsvideo{24@1}{7zI5zsPuRVs} % Projectors,bars and kets - Lec 04 - Frederic Schuller
\tsvideo{25@1}{6ad9V8gvyBQ} % Measure Theory  -Lec05- Frederic Schuller
\tsvideo{26@1}{ot253Lhx2_o} % Integration of measurable functions - Lec06 - Frederic Schuller
\tsvideo{27@1}{22P_hcZID8w} % Self adjoint and essentially self-adjoint operators - Lec 07 - Frederic Schuller
\tsvideo{28@1}{DWwOhIJQexY} % Spectra and perturbation theory - L08 - Frederic Schuller
\tsvideo{29@1}{ntyMY2oXz1k} % Case study: momentum operator - Lec09 - Frederic Schuller
\tsvideo{30@1}{IIF8XfMda4k} % Inverse Spectral Theorem - L10 - Frederic Schuller
\tsvideo{31@1}{sPHpioiCDoM} % Spectral Theorem - L11 - Frederic Schuller
\tsvideo{32@1}{tOp2rdvOmd0} % Stone's theorem & construction of observables - L12 - Frederic Schuller
\tsvideo{33@1}{-SM3qnK24BU} % Spin - L13 - Frederic Schuller
\tsvideo{34@1}{OZ1WCyJmjgo} % Composite systems - L14 - Frederic Schuller
\tsvideo{35@1}{gr0acx-dpPw} % Total spin of composite system - L15 - Frederic Schuller
\tsvideo{36@1}{UtO0qKmljW4} % Quantum Harmonic Oscillator - L16 - Frederic Schuller
\tsvideo{37@1}{s3I_MGfGm-w} % Quantum Harmonic Oscillator - L17 - Frederic Schuller
\tsvideo{38@1}{M-F70q13Jro} % The Fourier Operator - L18 - Frederic Schuller
\tsvideo{39@1}{RnPy8yYBRwg} % The Schrodinger Operator - L19 - Frederic Schuller
\tsvideo{40@1}{n2u41lXH80A} % Periodic potentials - L20 - Frederic Schuller
\tsvideo{41@1}{09GXy9ZH278} % Periodic potentials - L21 - Frederic Schuller

\makeatletter % consecutive tags may break between them, but never before punctuation
\DeclareRobustCommand{\ts}[3]{\mbox{\footnotesize\sffamily\color{srccol}\tslink{#1}{#2}{#3}{\color{srccol}[$\blacktriangleright$\,L#1\,@\,#3]}}\@ifnextchar\ts{\allowbreak}{}}
% link to a whole video (no time), e.g. in the Sources box: \seg{4}{2}{text}
\newcommand{\seg}[3]{\tslink{#1}{#2}{0}{#3}}
\makeatother

% ---------- maths shortcuts ----------
\newcommand{\R}{\mathbb{R}}
\newcommand{\N}{\mathbb{N}}
\newcommand{\E}{\mathbb{E}}
\newcommand{\Prob}{\mathbb{P}}
\newcommand{\Q}{\mathbb{Q}}
\DeclareMathOperator{\Var}{Var}
\DeclareMathOperator{\Cov}{Cov}
\DeclareMathOperator{\Corr}{Corr}
\DeclareMathOperator{\rank}{rank}
\DeclareMathOperator{\diag}{diag}
\newcommand{\C}{\mathbb{C}}
\newcommand{\K}{\mathcal{K}}  % Koopman operator
\newcommand{\vb}[1]{\mathbf{#1}} % vectors and matrices of data: \vb x, \vb X
\newcommand{\dd}{\,\mathrm{d}}
\newcommand{\ind}{\mathbf{1}}
\newcommand{\eqd}{\stackrel{d}{=}}

% ---------- code: the Matlab demos of the lectures, rewritten in Python ----------
%  \begin{code}[topic] ... \end{code}   (a listing in a grey-blue box; must be in the text column)
\lstdefinestyle{py}{language=Python,basicstyle=\ttfamily\footnotesize,columns=fullflexible,keepspaces,
  keywordstyle=\color{defcol}\bfseries,commentstyle=\color{intcol}\itshape,stringstyle=\color{fincol},
  showstringspaces=false,breaklines=true,upquote=true}
\newtcblisting{code}[1][]{notebox=srccol,colback=srccol!4,listing only,listing options={style=py},
  title={\boxtitle{Python}{#1}}}
%  \codefile{topic}{../code/chNN_name.py}: the same box, reading a script of ../code (run it from the
%  repository root to regenerate the figures in figures/), so notes and scripts never diverge
\newcommand{\codefile}[2]{\tcbinputlisting{notebox=srccol,colback=srccol!4,listing only,
  listing options={style=py},listing file={#2},title={\boxtitle{Python}{#1}~{\normalfont\ttfamily\footnotesize(\detokenize{#2})}}}}

% ---------- file references ----------
\DeclareRobustCommand{\file}[1]{\texttt{\small\detokenize{#1}}}


% Compile only some chapters while working on them, e.g.:
% \includeonly{chapters/ch01-<slug>}

\begin{document}

\frontmatter
\begin{titlepage}
  \centering
  \vspace*{3cm}
  {\Huge\bfseries Quantum Theory\par}
  \vspace{1cm}
  {\Large Lecture notes\par}
  \vspace{0.5cm}
  {\large based on the video lectures by Frederic P.~Schuller\\
  (Friedrich-Alexander-Universit\"at Erlangen-N\"urnberg)\\
  Part I: \emph{Theoretische Quantenmechanik} (TP-2, 20 lectures, German, winter semester 2014/15)\\Part II: \emph{Lectures on Quantum Theory} (21 lectures, English, summer semester 2015)\par}
  \vspace{2cm}
  {\large Personal study notes of Alessio Martini\par}
  \vfill
  \begin{minipage}{0.85\textwidth}\small
    These notes are a personal, non-commercial study aid derived from the recorded lectures by Frederic P.~Schuller
    (YouTube and \url{https://www.fau.tv/series/theoretische-physik-2-theoretische-quantenmechanik-tp-2-w14} and \url{https://www.fau.tv/series/lectures-on-quantum-theory-elite-graduate-programme-2015-measurements-s15}).
    Lecture content is paraphrased from the (automatic) transcripts of the videos, not transcribed;
    the blackboard notation is reconstructed from what is said, and doubtful points are marked~[?].
    Any mistakes are mine. All rights on the original lectures belong to their authors.
  \end{minipage}
  \vspace{1cm}
\end{titlepage}

\chapter*{How to use these notes}
\addcontentsline{toc}{chapter}{How to use these notes}
\begin{paracol}{2}

There is one chapter per lecture, and the sections of a chapter follow the lecture in the order in which
Schuller presents them. Every chapter opens with a \textbf{Sources} box (the video, its length, what it
contains) and a \textbf{Reading} box pointing to the matching chapters of the other Schuller notes sets.
Definitions, theorems and proofs are written as on his blackboard; everything else he says (motivation,
remarks, digressions, jokes with a point) is kept, the side remarks in the right column.

\paragraph{Boxes.}
\begin{itemize}
  \item \textcolor{defcol}{\textbf{Definitions}} (blue), \textcolor{thmcol}{\textbf{theorems, propositions, corollaries}} (red),
        \textcolor{exercol}{\textbf{exercises}} and \textcolor{srccol}{\textbf{remarks}} are marked with a
        coloured bar on the left. A definition box titled \emph{New concepts} marks the first time a concept appears.
  \item \textcolor{keycol}{\textbf{Key concepts}} at the end of each chapter: retrieval questions, then the concepts.
  \item \textcolor{intucol}{\textbf{Intuition}} boxes give the picture behind a definition.
  \item \textcolor{srccol}{\textbf{Reading}} boxes point to the other notes sets and to Schuller's papers.
  \item \textbf{Beyond the core} blocks (grey bar on the left, smaller type) hold material that can be
        skipped on a first reading.
\end{itemize}

\paragraph{The right column.} The page is laid out for a screen: the text is always on the left and a
column on the right holds \textcolor{intcol}{\textbf{examples}}, \textcolor{thmcol}{\textbf{proofs}},
\textcolor{srccol}{\textbf{course notes}}, \textcolor{histcol}{\textbf{history}} boxes and
\textcolor{histcol}{\textbf{asides}}, your \textcolor{notecol!80!black}{\textbf{notes}} with their
\textcolor{anscol}{\textbf{answers}}, and \textcolor{exercol}{\textbf{feedback}} on your solutions.

\paragraph{Video timestamps.} A tag such as \ts{1}{1}{28:22}\ means: lecture~1, at minute 28:22 of the
video. Every tag is a link to the YouTube video at that second.

\paragraph{Two courses.} Chapters 1--20 follow the German course TP-2 (written in English; the first occurrence of a term keeps the German in parentheses, e.g.\ phase space (Phasenraum)), chapters 21--41 follow the English course \emph{Lectures on Quantum Theory}. Topics that appear in both are cross-referenced, not merged.

\paragraph{Your annotations.} While reading, write directly in the chapter files. The macro
\verb|\note{...}| prints a comment in the right column next to the text where it is written. Requests
about the notes themselves are collected in the chapter \emph{Requests about these notes}.

\flushside
\end{paracol}

% !TEX root = ../main.tex
\chapter*{Requests about these notes}\label{ch:requests}
\addcontentsline{toc}{chapter}{Requests about these notes}
\markboth{Requests about these notes}{}
\begin{paracol}{2}

Your margin notes are of two kinds: notes on the \emph{content} (a doubt, an idea), which stay next to the
text they refer to, and requests about the \emph{notes themselves} (layout, boxes, conventions), which are
collected here with what was done about them.

\paragraph{How the margin notes are marked.}
\begin{itemize}
  \item \verb|\note{...}|: a \emph{new} note, not yet read. ``Read my notes'' means these, and only these.
  \item A note that has been answered becomes one box with the question on top and a dark yellow
        \textbf{ANSWER} bar below it (\verb|\begin{answer}{question}[topic]| \dots\ \verb|\end{answer}|).
  \item \verb|\note*{...}|: a request that has been carried out (green DONE tag).
  \item \verb|\note+{...}|: a comment that has been read and needs no answer (blue SEEN tag).
  \item \verb|\feedback{...}|: a comment on one of your own solutions (indigo).
\end{itemize}

\section*{The list}
No requests yet.

\flushside
\end{paracol}


\tableofcontents

\mainmatter
\mainmatter
% BEGIN CHAPTERS (tools/stubs.py)
\part{TP-2: Theoretische Quantenmechanik}
% !TEX root = ../main.tex
% STUB: generated by tools/stubs.py; write the chapter from the source (transcript in materials/transcripts/ if any)
\chapter{Experimentum crucis and quantum probabilities}\label{ch:c01}
\begin{paracol}{2}

\begin{sources}
\begin{tabularx}{\linewidth}{@{}l l L@{}}
\toprule
\textbf{Source} & \textbf{Length} & \textbf{Content}\\
\midrule
\seg{1}{1}{Experimentum crucis and quantum probabilities} & 1:49:00 & (to be written)\\
\bottomrule
\end{tabularx}

\smallskip
\textbf{Transcript} \file{materials/transcripts/C01-V01.txt}.
\end{sources}

\begin{reading}[Where this chapter sits in the other notes]
(to be written)
\end{reading}

\emph{This chapter has not been written yet.}

\flushside
\end{paracol}

% !TEX root = ../main.tex
% STUB: generated by tools/stubs.py; write the chapter from the source (transcript in materials/transcripts/ if any)
\chapter{Path integral and elementary amplitudes}\label{ch:c02}
\begin{paracol}{2}

\begin{sources}
\begin{tabularx}{\linewidth}{@{}l l L@{}}
\toprule
\textbf{Source} & \textbf{Length} & \textbf{Content}\\
\midrule
\seg{2}{1}{Path integral and elementary amplitudes} & 1:36:51 & (to be written)\\
\bottomrule
\end{tabularx}

\smallskip
\textbf{Transcript} \file{materials/transcripts/C02-V01.txt}.
\end{sources}

\begin{reading}[Where this chapter sits in the other notes]
(to be written)
\end{reading}

\emph{This chapter has not been written yet.}

\flushside
\end{paracol}

% !TEX root = ../main.tex
% STUB: generated by tools/stubs.py; write the chapter from the source (transcript in materials/transcripts/ if any)
\chapter{Wave functions and the Schrödinger equation}\label{ch:c03}
\begin{paracol}{2}

\begin{sources}
\begin{tabularx}{\linewidth}{@{}l l L@{}}
\toprule
\textbf{Source} & \textbf{Length} & \textbf{Content}\\
\midrule
\seg{3}{1}{Wave functions and the Schrödinger equation} & 1:48:30 & (to be written)\\
\bottomrule
\end{tabularx}

\smallskip
\textbf{Transcript} \file{materials/transcripts/C03-V01.txt}.
\end{sources}

\begin{reading}[Where this chapter sits in the other notes]
(to be written)
\end{reading}

\emph{This chapter has not been written yet.}

\flushside
\end{paracol}

% !TEX root = ../main.tex
% STUB: generated by tools/stubs.py; write the chapter from the source (transcript in materials/transcripts/ if any)
\chapter{Single slit and uncertainty relation}\label{ch:c04}
\begin{paracol}{2}

\begin{sources}
\begin{tabularx}{\linewidth}{@{}l l L@{}}
\toprule
\textbf{Source} & \textbf{Length} & \textbf{Content}\\
\midrule
\seg{4}{1}{Single slit and uncertainty relation} & 1:44:05 & (to be written)\\
\bottomrule
\end{tabularx}

\smallskip
\textbf{Transcript} \file{materials/transcripts/C04-V01.txt}.
\end{sources}

\begin{reading}[Where this chapter sits in the other notes]
(to be written)
\end{reading}

\emph{This chapter has not been written yet.}

\flushside
\end{paracol}

% !TEX root = ../main.tex
% STUB: generated by tools/stubs.py; write the chapter from the source (transcript in materials/transcripts/ if any)
\chapter{Hilbert spaces (Hilbert(t)räume)}\label{ch:c05}
\begin{paracol}{2}

\begin{sources}
\begin{tabularx}{\linewidth}{@{}l l L@{}}
\toprule
\textbf{Source} & \textbf{Length} & \textbf{Content}\\
\midrule
\seg{5}{1}{Hilbert spaces (Hilbert(t)räume)} & 1:48:15 & (to be written)\\
\bottomrule
\end{tabularx}

\smallskip
\textbf{Transcript} \file{materials/transcripts/C05-V01.txt}.
\end{sources}

\begin{reading}[Where this chapter sits in the other notes]
(to be written)
\end{reading}

\emph{This chapter has not been written yet.}

\flushside
\end{paracol}

% !TEX root = ../main.tex
% STUB: generated by tools/stubs.py; write the chapter from the source (transcript in materials/transcripts/ if any)
\chapter{Spectral theorem (toddler version)}\label{ch:c06}
\begin{paracol}{2}

\begin{sources}
\begin{tabularx}{\linewidth}{@{}l l L@{}}
\toprule
\textbf{Source} & \textbf{Length} & \textbf{Content}\\
\midrule
\seg{6}{1}{Spectral theorem (toddler version)} & 1:45:03 & (to be written)\\
\bottomrule
\end{tabularx}

\smallskip
\textbf{Transcript} \file{materials/transcripts/C06-V01.txt}.
\end{sources}

\begin{reading}[Where this chapter sits in the other notes]
(to be written)
\end{reading}

\emph{This chapter has not been written yet.}

\flushside
\end{paracol}

% !TEX root = ../main.tex
% STUB: generated by tools/stubs.py; write the chapter from the source (transcript in materials/transcripts/ if any)
\chapter{Spectral theorem (general)}\label{ch:c07}
\begin{paracol}{2}

\begin{sources}
\begin{tabularx}{\linewidth}{@{}l l L@{}}
\toprule
\textbf{Source} & \textbf{Length} & \textbf{Content}\\
\midrule
\seg{7}{1}{Spectral theorem (general)} & 1:44:55 & (to be written)\\
\bottomrule
\end{tabularx}

\smallskip
\textbf{Transcript} \file{materials/transcripts/C07-V01.txt}.
\end{sources}

\begin{reading}[Where this chapter sits in the other notes]
(to be written)
\end{reading}

\emph{This chapter has not been written yet.}

\flushside
\end{paracol}

% !TEX root = ../main.tex
% STUB: generated by tools/stubs.py; write the chapter from the source (transcript in materials/transcripts/ if any)
\chapter{Rigged Hilbert spaces (aufgetakelte Hilberträume)}\label{ch:c08}
\begin{paracol}{2}

\begin{sources}
\begin{tabularx}{\linewidth}{@{}l l L@{}}
\toprule
\textbf{Source} & \textbf{Length} & \textbf{Content}\\
\midrule
\seg{8}{1}{Rigged Hilbert spaces (aufgetakelte Hilberträume)} & 1:46:22 & (to be written)\\
\bottomrule
\end{tabularx}

\smallskip
\textbf{Transcript} \file{materials/transcripts/C08-V01.txt}.
\end{sources}

\begin{reading}[Where this chapter sits in the other notes]
(to be written)
\end{reading}

\emph{This chapter has not been written yet.}

\flushside
\end{paracol}

% !TEX root = ../main.tex
% STUB: generated by tools/stubs.py; write the chapter from the source (transcript in materials/transcripts/ if any)
\chapter{General framework of QM: states and observables}\label{ch:c09}
\begin{paracol}{2}

\begin{sources}
\begin{tabularx}{\linewidth}{@{}l l L@{}}
\toprule
\textbf{Source} & \textbf{Length} & \textbf{Content}\\
\midrule
\seg{9}{1}{General framework of QM: states and observables} & 1:49:12 & (to be written)\\
\bottomrule
\end{tabularx}

\smallskip
\textbf{Transcript} \file{materials/transcripts/C09-V01.txt}.
\end{sources}

\begin{reading}[Where this chapter sits in the other notes]
(to be written)
\end{reading}

\emph{This chapter has not been written yet.}

\flushside
\end{paracol}

% !TEX root = ../main.tex
% STUB: generated by tools/stubs.py; write the chapter from the source (transcript in materials/transcripts/ if any)
\chapter{General framework of QM: best possible prediction of measurement results}\label{ch:c10}
\begin{paracol}{2}

\begin{sources}
\begin{tabularx}{\linewidth}{@{}l l L@{}}
\toprule
\textbf{Source} & \textbf{Length} & \textbf{Content}\\
\midrule
\seg{10}{1}{General framework of QM: best possible prediction of measurement results} & 1:22:12 & (to be written)\\
\bottomrule
\end{tabularx}

\smallskip
\textbf{Transcript} \file{materials/transcripts/C10-V01.txt}.
\end{sources}

\begin{reading}[Where this chapter sits in the other notes]
(to be written)
\end{reading}

\emph{This chapter has not been written yet.}

\flushside
\end{paracol}

% !TEX root = ../main.tex
% STUB: generated by tools/stubs.py; write the chapter from the source (transcript in materials/transcripts/ if any)
\chapter{General framework of QM: dynamics at and between measurements}\label{ch:c11}
\begin{paracol}{2}

\begin{sources}
\begin{tabularx}{\linewidth}{@{}l l L@{}}
\toprule
\textbf{Source} & \textbf{Length} & \textbf{Content}\\
\midrule
\seg{11}{1}{General framework of QM: dynamics at and between measurements} & 1:48:48 & (to be written)\\
\bottomrule
\end{tabularx}

\smallskip
\textbf{Transcript} \file{materials/transcripts/C11-V01.txt}.
\end{sources}

\begin{reading}[Where this chapter sits in the other notes]
(to be written)
\end{reading}

\emph{This chapter has not been written yet.}

\flushside
\end{paracol}

% !TEX root = ../main.tex
% STUB: generated by tools/stubs.py; write the chapter from the source (transcript in materials/transcripts/ if any)
\chapter{From the experimentum crucis to the general framework of QM}\label{ch:c12}
\begin{paracol}{2}

\begin{sources}
\begin{tabularx}{\linewidth}{@{}l l L@{}}
\toprule
\textbf{Source} & \textbf{Length} & \textbf{Content}\\
\midrule
From the experimentum crucis to the general framework of QM & 1:24:17 & (to be written) (no YouTube id; FAU.tv: \url{https://www.fau.tv/series/theoretische-physik-2-theoretische-quantenmechanik-tp-2-w14/12-theoretische-quantenmechanik-vom-experimentum-crucis-zum-allg-rahmenwerk-der-qm-w14})\\
\bottomrule
\end{tabularx}

\smallskip
\textbf{Transcript} \file{materials/transcripts/C12-V01.txt}.
\end{sources}

\begin{reading}[Where this chapter sits in the other notes]
(to be written)
\end{reading}

\emph{This chapter has not been written yet.}

\flushside
\end{paracol}

% !TEX root = ../main.tex
% STUB: generated by tools/stubs.py; write the chapter from the source (transcript in materials/transcripts/ if any)
\chapter{Representations with respect to arbitrary observables}\label{ch:c13}
\begin{paracol}{2}

\begin{sources}
\begin{tabularx}{\linewidth}{@{}l l L@{}}
\toprule
\textbf{Source} & \textbf{Length} & \textbf{Content}\\
\midrule
Representations with respect to arbitrary observables & 1:18:33 & (to be written) (no YouTube id; FAU.tv: \url{https://www.fau.tv/series/theoretische-physik-2-theoretische-quantenmechanik-tp-2-w14/13-theoretische-quantenmechanik-darstellungen-bezuglich-beliebiger-observablen-w14})\\
\bottomrule
\end{tabularx}

\smallskip
\textbf{Transcript} \file{materials/transcripts/C13-V01.txt}.
\end{sources}

\begin{reading}[Where this chapter sits in the other notes]
(to be written)
\end{reading}

\emph{This chapter has not been written yet.}

\flushside
\end{paracol}

% !TEX root = ../main.tex
% STUB: generated by tools/stubs.py; write the chapter from the source (transcript in materials/transcripts/ if any)
\chapter{Symmetries, in particular rotation symmetry and the angular-momentum algebra}\label{ch:c14}
\begin{paracol}{2}

\begin{sources}
\begin{tabularx}{\linewidth}{@{}l l L@{}}
\toprule
\textbf{Source} & \textbf{Length} & \textbf{Content}\\
\midrule
Symmetries, in particular rotation symmetry and the angular-momentum algebra & 1:26:40 & (to be written) (no YouTube id; FAU.tv: \url{https://www.fau.tv/series/theoretische-physik-2-theoretische-quantenmechanik-tp-2-w14/14-theoretische-quantenmechanik-symmetrien-speziell-rotationssymmetrie-amp-drehimpulsalgebra-w14})\\
\bottomrule
\end{tabularx}

\smallskip
\textbf{Transcript} \file{materials/transcripts/C14-V01.txt}.
\end{sources}

\begin{reading}[Where this chapter sits in the other notes]
(to be written)
\end{reading}

\emph{This chapter has not been written yet.}

\flushside
\end{paracol}

% !TEX root = ../main.tex
% STUB: generated by tools/stubs.py; write the chapter from the source (transcript in materials/transcripts/ if any)
\chapter{"Gluing together" quantum systems}\label{ch:c15}
\begin{paracol}{2}

\begin{sources}
\begin{tabularx}{\linewidth}{@{}l l L@{}}
\toprule
\textbf{Source} & \textbf{Length} & \textbf{Content}\\
\midrule
"Gluing together" quantum systems & 1:44:04 & (to be written) (no YouTube id; FAU.tv: \url{https://www.fau.tv/series/theoretische-physik-2-theoretische-quantenmechanik-tp-2-w14/15-theoretische-quantenmechanik-quotzusammenklebenquot-von-quantenmechanischen-systemen-w14})\\
\bottomrule
\end{tabularx}

\smallskip
\textbf{Transcript} \file{materials/transcripts/C15-V01.txt}.
\end{sources}

\begin{reading}[Where this chapter sits in the other notes]
(to be written)
\end{reading}

\emph{This chapter has not been written yet.}

\flushside
\end{paracol}

% !TEX root = ../main.tex
% STUB: generated by tools/stubs.py; write the chapter from the source (transcript in materials/transcripts/ if any)
\chapter{Addition of angular momenta and Clebsch-Gordan coefficients}\label{ch:c16}
\begin{paracol}{2}

\begin{sources}
\begin{tabularx}{\linewidth}{@{}l l L@{}}
\toprule
\textbf{Source} & \textbf{Length} & \textbf{Content}\\
\midrule
Addition of angular momenta and Clebsch-Gordan coefficients & 1:46:30 & (to be written) (no YouTube id; FAU.tv: \url{https://www.fau.tv/series/theoretische-physik-2-theoretische-quantenmechanik-tp-2-w14/16-theoretische-quantenmechanik-addition-von-drehimpulsoperatoren-amp-clebsch-gordan-koeffizienten-w14})\\
\bottomrule
\end{tabularx}

\smallskip
\textbf{Transcript} \file{materials/transcripts/C16-V01.txt}.
\end{sources}

\begin{reading}[Where this chapter sits in the other notes]
(to be written)
\end{reading}

\emph{This chapter has not been written yet.}

\flushside
\end{paracol}

% !TEX root = ../main.tex
% STUB: generated by tools/stubs.py; write the chapter from the source (transcript in materials/transcripts/ if any)
\chapter{The quantum two-body problem}\label{ch:c17}
\begin{paracol}{2}

\begin{sources}
\begin{tabularx}{\linewidth}{@{}l l L@{}}
\toprule
\textbf{Source} & \textbf{Length} & \textbf{Content}\\
\midrule
The quantum two-body problem & 1:23:55 & (to be written) (no YouTube id; FAU.tv: \url{https://www.fau.tv/series/theoretische-physik-2-theoretische-quantenmechanik-tp-2-w14/17-theoretische-quantenmechanik-quantenmechanisches-zweikorperproblem-w14})\\
\bottomrule
\end{tabularx}

\smallskip
\textbf{Transcript} \file{materials/transcripts/C17-V01.txt}.
\end{sources}

\begin{reading}[Where this chapter sits in the other notes]
(to be written)
\end{reading}

\emph{This chapter has not been written yet.}

\flushside
\end{paracol}

% !TEX root = ../main.tex
% STUB: generated by tools/stubs.py; write the chapter from the source (transcript in materials/transcripts/ if any)
\chapter{The hydrogen atom}\label{ch:c18}
\begin{paracol}{2}

\begin{sources}
\begin{tabularx}{\linewidth}{@{}l l L@{}}
\toprule
\textbf{Source} & \textbf{Length} & \textbf{Content}\\
\midrule
The hydrogen atom & 1:05:31 & (to be written) (no YouTube id; FAU.tv: \url{https://www.fau.tv/series/theoretische-physik-2-theoretische-quantenmechanik-tp-2-w14/18-theoretische-quantenmechanik-das-wasserstoffatom-w14})\\
\bottomrule
\end{tabularx}

\smallskip
\textbf{Transcript} \file{materials/transcripts/C18-V01.txt}.
\end{sources}

\begin{reading}[Where this chapter sits in the other notes]
(to be written)
\end{reading}

\emph{This chapter has not been written yet.}

\flushside
\end{paracol}

% !TEX root = ../main.tex
% STUB: generated by tools/stubs.py; write the chapter from the source (transcript in materials/transcripts/ if any)
\chapter{The general uncertainty relation}\label{ch:c19}
\begin{paracol}{2}

\begin{sources}
\begin{tabularx}{\linewidth}{@{}l l L@{}}
\toprule
\textbf{Source} & \textbf{Length} & \textbf{Content}\\
\midrule
The general uncertainty relation & 1:28:11 & (to be written) (no YouTube id; FAU.tv: \url{https://www.fau.tv/series/theoretische-physik-2-theoretische-quantenmechanik-tp-2-w14/19-theoretische-quantenmechanik-allgemeine-unscharferlation-w14})\\
\bottomrule
\end{tabularx}

\smallskip
\textbf{Transcript} \file{materials/transcripts/C19-V01.txt}.
\end{sources}

\begin{reading}[Where this chapter sits in the other notes]
(to be written)
\end{reading}

\emph{This chapter has not been written yet.}

\flushside
\end{paracol}

% !TEX root = ../main.tex
% STUB: generated by tools/stubs.py; write the chapter from the source (transcript in materials/transcripts/ if any)
\chapter{One-electron model of solid-state physics}\label{ch:c20}
\begin{paracol}{2}

\begin{sources}
\begin{tabularx}{\linewidth}{@{}l l L@{}}
\toprule
\textbf{Source} & \textbf{Length} & \textbf{Content}\\
\midrule
One-electron model of solid-state physics & 1:21:06 & (to be written) (no YouTube id; FAU.tv: \url{https://www.fau.tv/series/theoretische-physik-2-theoretische-quantenmechanik-tp-2-w14/20-theoretische-quantenmechanik-ein-elektronen-modell-der-festkorperphysik-w14})\\
\bottomrule
\end{tabularx}

\smallskip
\textbf{Transcript} \file{materials/transcripts/C20-V01.txt}.
\end{sources}

\begin{reading}[Where this chapter sits in the other notes]
(to be written)
\end{reading}

\emph{This chapter has not been written yet.}

\flushside
\end{paracol}

\part{Lectures on Quantum Theory}
% !TEX root = ../main.tex
% STUB: generated by tools/stubs.py; write the chapter from the source (transcript in materials/transcripts/ if any)
\chapter{Axioms of Quantum Mechanics}\label{ch:c21}
\begin{paracol}{2}

\begin{sources}
\begin{tabularx}{\linewidth}{@{}l l L@{}}
\toprule
\textbf{Source} & \textbf{Length} & \textbf{Content}\\
\midrule
\seg{21}{1}{Axioms of Quantum Mechanics} & 2:09:35 & (to be written)\\
\bottomrule
\end{tabularx}

\smallskip
\textbf{Transcript} \file{materials/transcripts/C21-V01.txt}.
\end{sources}

\begin{reading}[Where this chapter sits in the other notes]
(to be written)
\end{reading}

\emph{This chapter has not been written yet.}

\flushside
\end{paracol}

% !TEX root = ../main.tex
% STUB: generated by tools/stubs.py; write the chapter from the source (transcript in materials/transcripts/ if any)
\chapter{Banach Spaces}\label{ch:c22}
\begin{paracol}{2}

\begin{sources}
\begin{tabularx}{\linewidth}{@{}l l L@{}}
\toprule
\textbf{Source} & \textbf{Length} & \textbf{Content}\\
\midrule
\seg{22}{1}{Banach Spaces} & 1:49:17 & (to be written)\\
\bottomrule
\end{tabularx}

\smallskip
\textbf{Transcript} \file{materials/transcripts/C22-V01.txt}.
\end{sources}

\begin{reading}[Where this chapter sits in the other notes]
(to be written)
\end{reading}

\emph{This chapter has not been written yet.}

\flushside
\end{paracol}

% !TEX root = ../main.tex
% STUB: generated by tools/stubs.py; write the chapter from the source (transcript in materials/transcripts/ if any)
\chapter{Separable Hilbert spaces}\label{ch:c23}
\begin{paracol}{2}

\begin{sources}
\begin{tabularx}{\linewidth}{@{}l l L@{}}
\toprule
\textbf{Source} & \textbf{Length} & \textbf{Content}\\
\midrule
\seg{23}{1}{Separable Hilbert spaces} & 1:48:28 & (to be written)\\
\bottomrule
\end{tabularx}

\smallskip
\textbf{Transcript} \file{materials/transcripts/C23-V01.txt}.
\end{sources}

\begin{reading}[Where this chapter sits in the other notes]
(to be written)
\end{reading}

\emph{This chapter has not been written yet.}

\flushside
\end{paracol}

% !TEX root = ../main.tex
% STUB: generated by tools/stubs.py; write the chapter from the source (transcript in materials/transcripts/ if any)
\chapter{Projectors,bars and kets}\label{ch:c24}
\begin{paracol}{2}

\begin{sources}
\begin{tabularx}{\linewidth}{@{}l l L@{}}
\toprule
\textbf{Source} & \textbf{Length} & \textbf{Content}\\
\midrule
\seg{24}{1}{Projectors,bars and kets} & 1:44:10 & (to be written)\\
\bottomrule
\end{tabularx}

\smallskip
\textbf{Transcript} \file{materials/transcripts/C24-V01.txt}.
\end{sources}

\begin{reading}[Where this chapter sits in the other notes]
(to be written)
\end{reading}

\emph{This chapter has not been written yet.}

\flushside
\end{paracol}

% !TEX root = ../main.tex
% STUB: generated by tools/stubs.py; write the chapter from the source (transcript in materials/transcripts/ if any)
\chapter{Measure Theory}\label{ch:c25}
\begin{paracol}{2}

\begin{sources}
\begin{tabularx}{\linewidth}{@{}l l L@{}}
\toprule
\textbf{Source} & \textbf{Length} & \textbf{Content}\\
\midrule
\seg{25}{1}{Measure Theory} & 1:45:50 & (to be written)\\
\bottomrule
\end{tabularx}

\smallskip
\textbf{Transcript} \file{materials/transcripts/C25-V01.txt}.
\end{sources}

\begin{reading}[Where this chapter sits in the other notes]
(to be written)
\end{reading}

\emph{This chapter has not been written yet.}

\flushside
\end{paracol}

% !TEX root = ../main.tex
% STUB: generated by tools/stubs.py; write the chapter from the source (transcript in materials/transcripts/ if any)
\chapter{Integration of measurable functions}\label{ch:c26}
\begin{paracol}{2}

\begin{sources}
\begin{tabularx}{\linewidth}{@{}l l L@{}}
\toprule
\textbf{Source} & \textbf{Length} & \textbf{Content}\\
\midrule
\seg{26}{1}{Integration of measurable functions} & 1:53:40 & (to be written)\\
\bottomrule
\end{tabularx}

\smallskip
\textbf{Transcript} \file{materials/transcripts/C26-V01.txt}.
\end{sources}

\begin{reading}[Where this chapter sits in the other notes]
(to be written)
\end{reading}

\emph{This chapter has not been written yet.}

\flushside
\end{paracol}

% !TEX root = ../main.tex
% STUB: generated by tools/stubs.py; write the chapter from the source (transcript in materials/transcripts/ if any)
\chapter{Self adjoint and essentially self-adjoint operators}\label{ch:c27}
\begin{paracol}{2}

\begin{sources}
\begin{tabularx}{\linewidth}{@{}l l L@{}}
\toprule
\textbf{Source} & \textbf{Length} & \textbf{Content}\\
\midrule
\seg{27}{1}{Self adjoint and essentially self-adjoint operators} & 1:42:29 & (to be written)\\
\bottomrule
\end{tabularx}

\smallskip
\textbf{Transcript} \file{materials/transcripts/C27-V01.txt}.
\end{sources}

\begin{reading}[Where this chapter sits in the other notes]
(to be written)
\end{reading}

\emph{This chapter has not been written yet.}

\flushside
\end{paracol}

% !TEX root = ../main.tex
% STUB: generated by tools/stubs.py; write the chapter from the source (transcript in materials/transcripts/ if any)
\chapter{Spectra and perturbation theory}\label{ch:c28}
\begin{paracol}{2}

\begin{sources}
\begin{tabularx}{\linewidth}{@{}l l L@{}}
\toprule
\textbf{Source} & \textbf{Length} & \textbf{Content}\\
\midrule
\seg{28}{1}{Spectra and perturbation theory} & 2:07:23 & (to be written)\\
\bottomrule
\end{tabularx}

\smallskip
\textbf{Transcript} \file{materials/transcripts/C28-V01.txt}.
\end{sources}

\begin{reading}[Where this chapter sits in the other notes]
(to be written)
\end{reading}

\emph{This chapter has not been written yet.}

\flushside
\end{paracol}

% !TEX root = ../main.tex
% STUB: generated by tools/stubs.py; write the chapter from the source (transcript in materials/transcripts/ if any)
\chapter{Case study: momentum operator}\label{ch:c29}
\begin{paracol}{2}

\begin{sources}
\begin{tabularx}{\linewidth}{@{}l l L@{}}
\toprule
\textbf{Source} & \textbf{Length} & \textbf{Content}\\
\midrule
\seg{29}{1}{Case study: momentum operator} & 1:50:15 & (to be written)\\
\bottomrule
\end{tabularx}

\smallskip
\textbf{Transcript} \file{materials/transcripts/C29-V01.txt}.
\end{sources}

\begin{reading}[Where this chapter sits in the other notes]
(to be written)
\end{reading}

\emph{This chapter has not been written yet.}

\flushside
\end{paracol}

% !TEX root = ../main.tex
% STUB: generated by tools/stubs.py; write the chapter from the source (transcript in materials/transcripts/ if any)
\chapter{Inverse Spectral Theorem}\label{ch:c30}
\begin{paracol}{2}

\begin{sources}
\begin{tabularx}{\linewidth}{@{}l l L@{}}
\toprule
\textbf{Source} & \textbf{Length} & \textbf{Content}\\
\midrule
\seg{30}{1}{Inverse Spectral Theorem} & 1:54:20 & (to be written)\\
\bottomrule
\end{tabularx}

\smallskip
\textbf{Transcript} \file{materials/transcripts/C30-V01.txt}.
\end{sources}

\begin{reading}[Where this chapter sits in the other notes]
(to be written)
\end{reading}

\emph{This chapter has not been written yet.}

\flushside
\end{paracol}

% !TEX root = ../main.tex
% STUB: generated by tools/stubs.py; write the chapter from the source (transcript in materials/transcripts/ if any)
\chapter{Spectral Theorem}\label{ch:c31}
\begin{paracol}{2}

\begin{sources}
\begin{tabularx}{\linewidth}{@{}l l L@{}}
\toprule
\textbf{Source} & \textbf{Length} & \textbf{Content}\\
\midrule
\seg{31}{1}{Spectral Theorem} & 1:59:41 & (to be written)\\
\bottomrule
\end{tabularx}

\smallskip
\textbf{Transcript} \file{materials/transcripts/C31-V01.txt}.
\end{sources}

\begin{reading}[Where this chapter sits in the other notes]
(to be written)
\end{reading}

\emph{This chapter has not been written yet.}

\flushside
\end{paracol}

% !TEX root = ../main.tex
% STUB: generated by tools/stubs.py; write the chapter from the source (transcript in materials/transcripts/ if any)
\chapter{Stone's theorem \& construction of observables}\label{ch:c32}
\begin{paracol}{2}

\begin{sources}
\begin{tabularx}{\linewidth}{@{}l l L@{}}
\toprule
\textbf{Source} & \textbf{Length} & \textbf{Content}\\
\midrule
\seg{32}{1}{Stone's theorem \& construction of observables} & 1:52:25 & (to be written)\\
\bottomrule
\end{tabularx}

\smallskip
\textbf{Transcript} \file{materials/transcripts/C32-V01.txt}.
\end{sources}

\begin{reading}[Where this chapter sits in the other notes]
(to be written)
\end{reading}

\emph{This chapter has not been written yet.}

\flushside
\end{paracol}

% !TEX root = ../main.tex
% STUB: generated by tools/stubs.py; write the chapter from the source (transcript in materials/transcripts/ if any)
\chapter{Spin}\label{ch:c33}
\begin{paracol}{2}

\begin{sources}
\begin{tabularx}{\linewidth}{@{}l l L@{}}
\toprule
\textbf{Source} & \textbf{Length} & \textbf{Content}\\
\midrule
\seg{33}{1}{Spin} & 2:04:23 & (to be written)\\
\bottomrule
\end{tabularx}

\smallskip
\textbf{Transcript} \file{materials/transcripts/C33-V01.txt}.
\end{sources}

\begin{reading}[Where this chapter sits in the other notes]
(to be written)
\end{reading}

\emph{This chapter has not been written yet.}

\flushside
\end{paracol}

% !TEX root = ../main.tex
% STUB: generated by tools/stubs.py; write the chapter from the source (transcript in materials/transcripts/ if any)
\chapter{Composite systems}\label{ch:c34}
\begin{paracol}{2}

\begin{sources}
\begin{tabularx}{\linewidth}{@{}l l L@{}}
\toprule
\textbf{Source} & \textbf{Length} & \textbf{Content}\\
\midrule
\seg{34}{1}{Composite systems} & 1:44:36 & (to be written)\\
\bottomrule
\end{tabularx}

\smallskip
\textbf{Transcript} \file{materials/transcripts/C34-V01.txt}.
\end{sources}

\begin{reading}[Where this chapter sits in the other notes]
(to be written)
\end{reading}

\emph{This chapter has not been written yet.}

\flushside
\end{paracol}

% !TEX root = ../main.tex
% STUB: generated by tools/stubs.py; write the chapter from the source (transcript in materials/transcripts/ if any)
\chapter{Total spin of composite system}\label{ch:c35}
\begin{paracol}{2}

\begin{sources}
\begin{tabularx}{\linewidth}{@{}l l L@{}}
\toprule
\textbf{Source} & \textbf{Length} & \textbf{Content}\\
\midrule
\seg{35}{1}{Total spin of composite system} & 1:55:14 & (to be written)\\
\bottomrule
\end{tabularx}

\smallskip
\textbf{Transcript} \file{materials/transcripts/C35-V01.txt}.
\end{sources}

\begin{reading}[Where this chapter sits in the other notes]
(to be written)
\end{reading}

\emph{This chapter has not been written yet.}

\flushside
\end{paracol}

% !TEX root = ../main.tex
% STUB: generated by tools/stubs.py; write the chapter from the source (transcript in materials/transcripts/ if any)
\chapter{Quantum Harmonic Oscillator}\label{ch:c36}
\begin{paracol}{2}

\begin{sources}
\begin{tabularx}{\linewidth}{@{}l l L@{}}
\toprule
\textbf{Source} & \textbf{Length} & \textbf{Content}\\
\midrule
\seg{36}{1}{Quantum Harmonic Oscillator} & 1:50:15 & (to be written)\\
\bottomrule
\end{tabularx}

\smallskip
\textbf{Transcript} \file{materials/transcripts/C36-V01.txt}.
\end{sources}

\begin{reading}[Where this chapter sits in the other notes]
(to be written)
\end{reading}

\emph{This chapter has not been written yet.}

\flushside
\end{paracol}

% !TEX root = ../main.tex
% STUB: generated by tools/stubs.py; write the chapter from the source (transcript in materials/transcripts/ if any)
\chapter{Quantum Harmonic Oscillator}\label{ch:c37}
\begin{paracol}{2}

\begin{sources}
\begin{tabularx}{\linewidth}{@{}l l L@{}}
\toprule
\textbf{Source} & \textbf{Length} & \textbf{Content}\\
\midrule
\seg{37}{1}{Quantum Harmonic Oscillator} & 1:53:33 & (to be written)\\
\bottomrule
\end{tabularx}

\smallskip
\textbf{Transcript} \file{materials/transcripts/C37-V01.txt}.
\end{sources}

\begin{reading}[Where this chapter sits in the other notes]
(to be written)
\end{reading}

\emph{This chapter has not been written yet.}

\flushside
\end{paracol}

% !TEX root = ../main.tex
% STUB: generated by tools/stubs.py; write the chapter from the source (transcript in materials/transcripts/ if any)
\chapter{The Fourier Operator}\label{ch:c38}
\begin{paracol}{2}

\begin{sources}
\begin{tabularx}{\linewidth}{@{}l l L@{}}
\toprule
\textbf{Source} & \textbf{Length} & \textbf{Content}\\
\midrule
\seg{38}{1}{The Fourier Operator} & 1:38:38 & (to be written)\\
\bottomrule
\end{tabularx}

\smallskip
\textbf{Transcript} \file{materials/transcripts/C38-V01.txt}.
\end{sources}

\begin{reading}[Where this chapter sits in the other notes]
(to be written)
\end{reading}

\emph{This chapter has not been written yet.}

\flushside
\end{paracol}

% !TEX root = ../main.tex
% STUB: generated by tools/stubs.py; write the chapter from the source (transcript in materials/transcripts/ if any)
\chapter{The Schrodinger Operator}\label{ch:c39}
\begin{paracol}{2}

\begin{sources}
\begin{tabularx}{\linewidth}{@{}l l L@{}}
\toprule
\textbf{Source} & \textbf{Length} & \textbf{Content}\\
\midrule
\seg{39}{1}{The Schrodinger Operator} & 1:52:11 & (to be written)\\
\bottomrule
\end{tabularx}

\smallskip
\textbf{Transcript} \file{materials/transcripts/C39-V01.txt}.
\end{sources}

\begin{reading}[Where this chapter sits in the other notes]
(to be written)
\end{reading}

\emph{This chapter has not been written yet.}

\flushside
\end{paracol}

% !TEX root = ../main.tex
% STUB: generated by tools/stubs.py; write the chapter from the source (transcript in materials/transcripts/ if any)
\chapter{Periodic potentials}\label{ch:c40}
\begin{paracol}{2}

\begin{sources}
\begin{tabularx}{\linewidth}{@{}l l L@{}}
\toprule
\textbf{Source} & \textbf{Length} & \textbf{Content}\\
\midrule
\seg{40}{1}{Periodic potentials} & 1:42:54 & (to be written)\\
\bottomrule
\end{tabularx}

\smallskip
\textbf{Transcript} \file{materials/transcripts/C40-V01.txt}.
\end{sources}

\begin{reading}[Where this chapter sits in the other notes]
(to be written)
\end{reading}

\emph{This chapter has not been written yet.}

\flushside
\end{paracol}

% !TEX root = ../main.tex
% STUB: generated by tools/stubs.py; write the chapter from the source (transcript in materials/transcripts/ if any)
\chapter{Periodic potentials}\label{ch:c41}
\begin{paracol}{2}

\begin{sources}
\begin{tabularx}{\linewidth}{@{}l l L@{}}
\toprule
\textbf{Source} & \textbf{Length} & \textbf{Content}\\
\midrule
\seg{41}{1}{Periodic potentials} & 1:33:45 & (to be written)\\
\bottomrule
\end{tabularx}

\smallskip
\textbf{Transcript} \file{materials/transcripts/C41-V01.txt}.
\end{sources}

\begin{reading}[Where this chapter sits in the other notes]
(to be written)
\end{reading}

\emph{This chapter has not been written yet.}

\flushside
\end{paracol}

% END CHAPTERS

\appendix
% !TEX root = ../main.tex
\chapter{Pseudo-exam and second reading}\label{app:exam}
\begin{paracol}{2}

The long exercises of each chapter, moved here verbatim so that a first reading of the chapters is not held
up by them. Use this appendix as a second reading or, once the course is finished, as a cumulative
self-test: retrieval spaced weeks after first study is among the best-supported ways to make material stick
(Cepeda et al.\ 2006; Adesope et al.\ 2017).

\flushside
\end{paracol}


\end{document}
